\chapter{The link to household surveys}\label{chap:microsim}

The model has a few household groups; poverty and inequality are measured on
household surveys. \manref{cge_microsim} passes the results of a simulation to
the survey data in memory, top-down and without feedback
\citep{bourguignon2008}. Each survey household $i$ belongs to a model
household group $h(i)$ and keeps its own income composition and, if given, its
own budget shares.

\section{Income}

For each group $h$ the model gives the ratio of new to base values of each
income source $k$; a zero base value is replaced by the ratio of the aggregate
shown:

\begin{cgetable}{p{0.30\linewidth}p{0.36\linewidth}p{0.28\linewidth}}
\toprule
Source & Ratio $\rho_{hk}$ & If the group's base value is zero\\
\midrule\endhead
\bottomrule\endlastfoot
factor $f$ & $\V{YIF}_{hf} / \V{YIF}^0_{hf}$ & $\V{YF}_f / \V{YF}^0_f$\\
government transfers (\cmd{gov}) & $\V{trnsfr}_{h,gov}\,\V{CPI} / (\V{trnsfr}^0_{h,gov}\,\V{CPI}^0)$ & $\V{CPI} / \V{CPI}^0$\\
transfers from abroad (\cmd{row}) & $\V{trnsfr}_{h,row}\,\V{EXR} / (\V{trnsfr}^0_{h,row}\,\V{EXR}^0)$ & $\V{EXR} / \V{EXR}^0$\\
transfers from other domestic institutions (\cmd{dom}) & $\sum_d \V{shii}_{hd}\,\V{DY}_d \big/ \sum_d \V{shii}_{hd}\,\V{DY}^0_d$ & $\V{CPI} / \V{CPI}^0$\\
indexed to prices (\cmd{cpi}) & $\V{CPI} / \V{CPI}^0$ & \\
unchanged (\cmd{fixed}) & 1 & \\
\end{cgetable}

Here $\V{DY}_d = (1 - \V{MPS}_d)(1 - \V{TINS}_d)\,\V{YI}_d$ is the income of
institution $d$ after direct taxes and savings. With income components
$y_{ik}$, the household's nominal consumption becomes
\begin{equation*}
C^1_i = C^0_i\,\frac{\sum_k \rho_{h(i)k}\,y_{ik}}{\sum_k y_{ik}}\;\kappa_{h(i)},
\qquad
\kappa_h = \frac{(1 - \V{MPS}_h)(1 - \V{TINS}_h)}{(1 - \V{MPS}^0_h)(1 - \V{TINS}^0_h)},
\end{equation*}
so that it follows its own income and its group's direct tax and savings
rates. A household with no income in the components given, or every household
when no components are given, follows its group's consumption:
$C^1_i = C^0_i\,\V{EH}_h / \V{EH}^0_h$.

\section{Prices and real consumption}

The household's price index is
\begin{equation*}
P_i = \sum_c s_{ic}\,\frac{\V{PH}_c}{\V{PH}^0_c},
\end{equation*}
at the prices that households pay, including VAT, with $s_{ic}$ its survey
budget shares when budget items are mapped to commodities, and the group's
base shares $\V{PH}^0_c\,\V{QH}^0_{ch} / \V{EH}^0_h$ otherwise. Home
consumption (\secref{sec:home}) enters the group's shares at the activity
price $\V{PA}_a$. Real consumption is $R_i = C^1_i / P_i$; for households with LES
demand this is a first-order (Laspeyres) approximation.

\section{Poverty and inequality}

With household size $n_i$ and survey weight $w_i$, poverty measures use
consumption per person, $r_i = R_i / n_i$ (base $C^0_i / n_i$), and person
weights $w_i\,n_i$. For a poverty line $z$ per person, the
Foster--Greer--Thorbecke measures are
\begin{equation*}
\V{FGT}_\alpha = \frac{\sum_i w_i\,n_i\,\mathbf{1}[r_i < z]\,\bigl((z - r_i)/z\bigr)^\alpha}{\sum_i w_i\,n_i},
\qquad \alpha = 0, 1, 2,
\end{equation*}
the headcount ratio, the poverty gap and the squared poverty gap. The Gini
coefficient is that of $r_i$ over persons. Both are computed before and after,
overall and by group.

\section{Distribution and incidence}

From version 0.8.0, a poverty line may differ by household, $z_i$, for
regional lines or lines that depend on household composition, and several
lines may be compared at once; for each, the number of poor,
$\sum_i w_i n_i\,\mathbf{1}[r_i < z_i]$, is given with the FGT measures.
Beside the Gini coefficient, the generalised entropy measures GE(0), the mean
log deviation, GE(1), the Theil index, and GE(2), the Atkinson indexes with
inequality aversion 0.5, 1 and 2, and the Palma and S80/S20 ratios describe
the distribution of consumption per person over persons
(\manref{cge_microsim} gives the formulas).

Incidence follows the same persons: ranked by base consumption per person
into quantile groups, deciles by default, each group's mean consumption
before and after, its shares of total consumption, and its share of the total
change show who gains and who loses. The growth incidence curve of
\citet{ravallion2003} compares the distributions instead: the growth of each
percentile of consumption per person, whoever holds it. Under the top-down
link, a household's rank can change only through its own income composition
and budget shares, so the two views agree closely unless the shock changes
relative incomes strongly.

Several simulations, such as alternative policies or the periods of a
dynamic run, can be passed to the survey in one call, with every result in one
dataset for tables and graphs.

\section{Welfare and the channels of change}

From version 0.8.1, each household's equivalent and compensating variations
measure its welfare change in money: with the expenditure function of its
group, the group's subsistence quantities scaled to the household, or, when
its budget is given, Cobb--Douglas preferences in its own budget shares. A
household that follows its group has its group's welfare change, as
\manref{cge_report} reports it, in proportion to its consumption, so the
survey and the model agree for the average household of each group.

A change in poverty or inequality has several causes: wages, profits,
transfers and remittances change by different amounts, the direct tax and
savings rates of the groups change, and prices change. Since a household's
consumption is the product of these changes, the contribution of each can be
found by switching the channels on one at a time; the order matters, and the
Shapley decomposition \citep{shorrocks2013} averages over every order, so that
the contributions add up exactly to the change. The change in poverty can
also be divided into growth, the change in the mean with an unchanged
distribution, and redistribution \citep{datt1992}.

\section{Fiscal incidence}

From version 0.8.2, the survey's own taxes and benefits can be analysed as
the Commitment to Equity method does \citep{lustig2018}: market income less
direct taxes, plus direct transfers, less indirect taxes, plus subsidies and
in-kind benefits, gives net market, disposable, consumable and final income,
and each instrument is described by its concentration in market income, its
progressivity (Kakwani index), coverage and marginal contribution to poverty
and inequality. What the model adds is the simulation: each instrument
changes as the model says that its base changes (direct taxes with incomes
and tax rates, indirect taxes with spending and the tax share of prices,
transfers as the government's transfers, in-kind benefits as the government's
or the health insurance fund's spending), so that fiscal incidence can be
compared before and after a policy, in a general equilibrium. In-kind health
benefits valued at cost, allocated to the households that use the services,
then follow the model's health spending: the incidence of a rise in the
health budget, from the government's purchases to the households that
benefit, is traced in one step.

\section{Graphs}

From version 0.9.6, \manref{cge_microsim} draws the results it computes, in
the style of figures in economics journals: the growth incidence curve,
incidence by quantile group and by household group, the change in poverty at
each line, the channels of the change (\cmd{decompose}), inequality by income
concept (\cmd{market()}) and the range across the draws of the elasticities
(\cmd{bands()}). The graphs plot the results of the tables above, and can be
saved as PDF or EPS files for publication.

\section{Weights, sampling error and the uncertainty of the model}

From version 0.8.3, three tools address what the top-down link leaves out.
First, a shock that moves workers between sectors changes who works where,
which the survey's fixed weights cannot show; reweighting the survey, so that
its totals of workers by sector (or any other characteristic) follow the
model, gives the households of the expanding sectors more weight
\citep{corral2022}. The new weights are as close as possible to the
survey's, in relative entropy. Second, the survey is a sample: with its
design declared by \cmd{svyset}, the mean and the poverty measures, and
their changes, have standard errors. Third, the model's elasticities are
uncertain: \manref{cge_sensitivity} draws them, keeps each draw as a
simulation, and the survey link gives the range of poverty and inequality
across the draws.

Survey households keep their own income composition and budget shares, so
their average change can differ from their group's: a household that lives on
wages gains more than its group when wages rise relative to other income. The
link is only as good as the mapping of income sources, which should cover
every substantial source, and the survey and the SAM should share a base year
where possible.
